\documentclass[11pt]{article}
\usepackage[T1]{fontenc}
\usepackage[margin=1in]{geometry}
\usepackage{amsmath,amssymb,amsthm,booktabs,graphicx,microtype}
\usepackage[hidelinks]{hyperref}
\newtheorem{theorem}{Theorem}[section]
\newtheorem{proposition}[theorem]{Proposition}
\newtheorem{lemma}[theorem]{Lemma}
\newtheorem{corollary}[theorem]{Corollary}
\theoremstyle{definition}
\newtheorem{definition}[theorem]{Definition}
\newtheorem{remark}[theorem]{Remark}
\newcommand{\Q}{\mathbb Q}
\newcommand{\C}{\mathbb C}
\newcommand{\E}{\mathbb E}
\newcommand{\GL}{\operatorname{GL}}
\newcommand{\SL}{\operatorname{SL}}
\newcommand{\Sym}{\operatorname{Sym}}
\newcommand{\Tr}{\operatorname{Tr}}
\newcommand{\Aut}{\operatorname{Aut}}
\newcommand{\dd}{\,\mathrm d}
\numberwithin{equation}{section}
\title{A Quartic Period Torsor and Optimal Positive Moment Bounds}
\author{Alexander Kolpakov\\[0.5ex]
\small Center for Science, Technology, Engineering, and Mathematics\\
\small University of Austin, Austin, Texas, United States\\
\small \texttt{akolpakov@uaustin.org}\\[0.75ex]
\small The Wolfram Institute for Computational Foundations of Science\\
\small \texttt{akolpakov@wolframinstitute.org}}
\date{Research manuscript, 20 October 2026}
\begin{document}
\maketitle
\begin{abstract}
The stable quartic integral has a rank-two differential period torsor.
We describe its local realizations, their combinatorial expansions, and the
additional data needed to recover the real integral. Every finite Hermitian
matrix size is a matrix coefficient of a symmetric power of the same system.
To compare approximations, we classify positive generalized Stieltjes
representations: a positive homogeneous interaction of degree $2m$ in $d$
Gaussian variables has exact order $d/(2m)$. The smallest admissible order
gives the strongest Gauss and Gauss--Radau bounds from a fixed number of
perturbative coefficients. For quartic and sextic interactions these bounds
converge at every positive coupling. In particular, they give certified
approximations to all scalar sextic correlation functions. At degree eight
and above, a moment-indeterminacy example shows why positivity and all
perturbative coefficients can still be insufficient. Exact computations
illustrate the bounds for Hermitian matrices of sizes one through four.
Worked examples on finite rewrite graphs and hypergraph incidence graphs
express nonlinear field correlations through the same quartic periods
and spanning-forest counts.
\end{abstract}

\section{Introduction}
An integral may have several combinatorial expansions with very different
numerical behavior. A useful common object should explain their relations,
specify the information each expansion retains, and help decide which one
to use. A torsor of period comparisons suggests such an object, but the
integral value itself is not a torsor. Even in a simple example one must
distinguish the differential equation, a basis of solutions, an integration
cycle, and an approximation to the chosen period.

We work first with
\begin{equation}\label{eq:scalar}
 q_j(t)=\frac1{\sqrt{2\pi}}\int_{\mathbb R}
 x^j\exp\left(-\frac{x^2}{2}-\frac{t x^4}{4}\right)\dd x,
 \qquad t>0.
\end{equation}
The two moments $q_0,q_2$ satisfy a rational differential system. Its
horizontal frames form a differential torsor. The Gaussian expansion at
$t=0$, the quartic expansion at infinity, and Taylor expansions at ordinary
points are local realizations of this system. The formal local
classification is complete in this specified category; it is not a
classification of arbitrary differential extension fields.

There is also a second, useful classification. Suppose that the Wick
coefficients are used to reconstruct a positive integral of the form
$\int(1+\lambda s)^{-\beta}\dd\rho(s)$. Which orders $\beta$ are possible,
and which gives the best bounds? For the normalized Hermitian quartic
partition function of size $M$, the answer is $\beta\ge M^2/4$.
The minimal order gives stronger moment bounds than every larger order,
at the same number of input coefficients and at every positive coupling.
Here ``best'' is an extremal statement within a stated class of positive
representations, rather than a claim about all numerical algorithms.

The ingredients have substantial precedents. Picard--Vessiot torsors and
formal local differential equations are standard \cite{Singer}.
Comparison torsors for periods are treated by Brown \cite{Brown};
exponential integrals belong to the framework of Fres\'an and Jossen
\cite{FresanJossen}. Wick expansions and moment determinants are standard
matrix-model tools \cite{EKR}, and the quartic weight has a well-developed
orthogonal-polynomial and Painlev\'e theory \cite{ClarksonJordaan}.
Generalized positive-kernel approximants already appear in Savit's
work on lattice theories \cite[Appendix B]{Savit}.
Generalized Stieltjes classes and their exact orders are studied by Sokal
and by Karp--Prilepkina \cite{Sokal,KarpPrilepkina}; positive quadrature
bounds are classical \cite{Gautschi,LRW}. We give elementary proofs of the
specific representation and order comparison used here.

Fiol, Gij\'on and Lejarza Alonso recently proved convergence of rational
Schwinger--Dyson approximants for all scalar quartic correlation functions
\cite{FGL}. Their introduction asks about other even potentials and about
products of Pad\'e approximants. Our sextic corollary addresses convergence
by a different construction: positive quadrature of reweighted Wick
moments. It does not identify the truncations of the sextic
Schwinger--Dyson hierarchy or prove that they are products of Pad\'e
approximants.

The graph-theoretic motivation comes from the local Whitehead chain and
its stationary measure, weighted by inverse automorphism orders \cite{Kolpakov}.
The present paper keeps that motivation separate from the analysis of
the integral. In particular, a restriction to some graphs is not
automatically a specialization of a period torsor, nor does it automatically
bound the error in an integral.
Section~\ref{sec:rewrite} gives a concrete application to finite graphs
of rewrite histories and to hypergraph incidence graphs.

\section{The differential torsor}\label{sec:torsor}
For a rapid-decay contour $\gamma$, put
$m_j^\gamma=\int_\gamma x^j e^{-x^2/2-tx^4/4}\dd x$.
Integration by parts and differentiation give
\begin{equation}\label{eq:ibp}
 t m_{j+4}=(j+1)m_j-m_{j+2},\qquad
 (m_j)'=-\tfrac14m_{j+4}.
\end{equation}
Thus the even moments form a rank-two system
\begin{equation}\label{eq:connection}
 v'=A(t)v,\qquad
 A(t)=\begin{pmatrix}-1/(4t)&1/(4t)\\
 1/(4t^2)&-(1+3t)/(4t^2)\end{pmatrix},\qquad
 v=\begin{pmatrix}m_0\\m_2\end{pmatrix}.
\end{equation}
Equivalently,
\begin{equation}\label{eq:scalarODE}
 16t^2 y''+(32t+4)y'+3y=0,\qquad m_2=4t y'+y.
\end{equation}
The full quartic twisted de Rham system also has an odd-moment part.
Throughout, the torsor refers to the even-moment subsystem
\eqref{eq:connection}, generated by $\dd x,x^2\dd x$.

Let $K=\overline{\Q}(t)$, with derivation $D=\mathrm d/\mathrm dt$.
The coordinate algebra
\begin{equation}\label{eq:torsoralgebra}
 \mathcal R=K[y_{11},y_{12},y_{21},y_{22},\det(Y)^{-1}],
 \qquad DY=AY,
\end{equation}
is the frame torsor with its differential structure. Explicitly, for a
differential extension $F/K$, its horizontal points are
\[
 \mathcal T^D(F)=\{Y\in\GL_2(F):DY=AY\}.
\]
If this set has a point, it has a free transitive right action of
$\GL_2(F^D)$. Indeed,
$D(Y_1^{-1}Y_2)=0$ for any two horizontal frames. This is the torsor
property: two frames differ by one uniquely determined constant matrix.
The underlying algebraic $\GL_2$-torsor is trivial. The differential
structure and the selected period point carry the information.

\begin{proposition}\label{prop:galois}
The differential Galois group of \eqref{eq:connection} over $K$ is
$\GL_2$. In particular, \eqref{eq:torsoralgebra} is its Picard--Vessiot
torsor algebra.
\end{proposition}
\begin{proof}
Set $z=1/(8t)$ and $y=t^{-1/2}e^z f(z)$. Equation
\eqref{eq:scalarODE} becomes the modified Bessel equation of order
$1/4$. The Bessel group at this order is $\SL_2$
\cite[Example 1.3.32]{Singer}; replacing its variable by $iz$ gives the
modified equation. The algebraic gauge and exponential scalar twist do
not change its projective group, which is consequently $\operatorname{PGL}_2$.
On the other hand,
\begin{equation}\label{eq:wronskian}
 D\log\det Y=\Tr A=-\frac1t-\frac1{4t^2},\qquad
 \det Y=c\,t^{-1}e^{1/(4t)}.
\end{equation}
This determinant has group $\mathbb G_m$: no nonzero power of the
exponential is algebraic over $K$. The identity component of the full
group has projective image $\operatorname{PGL}_2$. A connected
three-dimensional subgroup with that projective image would have finite
scalar kernel, and hence no nontrivial character to $\mathbb G_m$.
The determinant image rules this out. The group therefore has dimension
four and equals $\GL_2$. The Picard--Vessiot torsor theorem
\cite[Proposition 1.3.20]{Singer} now identifies its algebra with
\eqref{eq:torsoralgebra}.
\end{proof}

For $t>0$, the contours $\mathbb R$ and $i\mathbb R$, oriented upwards,
give a period frame
\begin{equation}\label{eq:periodmatrix}
 P=\begin{pmatrix}m_0^{\mathbb R}&m_0^{i\mathbb R}\\
 m_2^{\mathbb R}&m_2^{i\mathbb R}\end{pmatrix}.
\end{equation}
These two columns will be seen to be independent below. They are points
over a complex analytic differential field; the gamma values in their
comparison matrix require extension of the constant field.
This differential torsor is not being identified with a full arithmetic
motivic comparison torsor at a fixed value of $t$. That would require
additional realization data and an arithmetic construction.

\section{Classification of the local realizations}\label{sec:local}
We allow completions at places of $K$, finite algebraic ramification,
and formal changes of basis. We also record when an exponential extension
is necessary. These operations give the following exhaustive list in
this local category.

\begin{theorem}\label{thm:local}
For \eqref{eq:connection} the local possibilities are as follows.
\begin{enumerate}
\item Every finite $a\ne0$ is ordinary. There is a unique horizontal frame
$U_a=I+O(t-a)$ over $\overline{\Q}[[t-a]]$.
\item At $t=0$, the formal scalar solutions have exponential factors
$1$ and $e^{1/(4t)}$. A full frame exists over
$\overline{\Q}((t))(e^{1/(4t)})$, but over no finite algebraic extension
of $\overline{\Q}((t))$.
\item At $t=\infty$, the equation is regular singular with scalar
exponents $1/4,3/4$ in $w=1/t$. A full frame exists over
$\overline{\Q}((w^{1/4}))$. In a cover $w=s^e$, it exists over
$\overline{\Q}((s))$ if and only if $4\mid e$.
\end{enumerate}
The analytic problem at zero also needs sector and Stokes information.
\end{theorem}
\begin{proof}
The first assertion follows by substituting $U_a=\sum_{n\ge0}U_n(t-a)^n$
in \eqref{eq:connection}. If $A=\sum A_j(t-a)^j$, then
\begin{equation}\label{eq:Taylor}
 (n+1)U_{n+1}=\sum_{j=0}^n A_jU_{n-j},\qquad U_0=I.
\end{equation}
At zero put
\begin{equation}\label{eq:weak}
 \widehat q(t)=\sum_{n\ge0}(-t)^n a_n,\qquad
 a_n=\frac{(4n-1)!!}{4^n n!}
     =\frac{4^n(1/4)_n(3/4)_n}{n!}.
\end{equation}
Direct substitution gives the two scalar solutions
$\widehat q(t)$ and $e^{1/(4t)}\widehat q(-t)$.
Applying $4tD+1$ to form their second components gives a frame of
determinant $-t^{-1}e^{1/(4t)}$. Its exponential determinant cannot
belong to any algebraic Laurent extension: its logarithmic derivative
has a double pole, whereas that of an algebraic Puiseux series has at
most a simple pole. This also proves the claimed obstruction.

At infinity, substituting $y=w^\rho(1+O(w))$ in
\eqref{eq:scalarODE} gives $(4\rho-1)(4\rho-3)=0$.
The exponents differ by a noninteger, so there are two logarithm-free
Frobenius series. Their fractional parts require ramification divisible
by four. This proves necessity and sufficiency. The explicit frame below
also proves sufficiency. These are instances of the usual formal local
classification \cite[Section 1.4]{Singer}.
\end{proof}

The coefficients in \eqref{eq:weak} count pairings of $4n$ Gaussian
half-edges, divided by the quartic vertex and exponential factors.
They do not select all contours. For the real contour the series is
asymptotic, and Taylor's integral remainder gives
\begin{equation}\label{eq:weakbound}
 \left|q_0(t)-\sum_{n=0}^K(-t)^n a_n\right|
 \le t^{K+1}a_{K+1}.
\end{equation}
Even partial sums are upper bounds, and odd partial sums are lower bounds.
The terms eventually increase at every $t>0$.

At infinity, scaling $x=t^{-1/4}y$ gives the convergent series
\begin{align}\label{eq:strong}
 q_0(t)&=\frac{t^{-1/4}}{2\sqrt\pi}
 \sum_{n\ge0}\frac{(-t^{-1/2})^n}{n!}\Gamma(n/2+1/4),\\
 q_2(t)&=\frac{t^{-3/4}}{\sqrt\pi}
 \sum_{n\ge0}\frac{(-t^{-1/2})^n}{n!}\Gamma(n/2+3/4).
 \nonumber
\end{align}
The sums are entire in $t^{-1/2}$. Dominated convergence follows from
$e^{-y^4/4+C y^2}$ on each compact parameter set. Truncating either
sum after degree $K$ has absolute remainder at most its positive
degree-$(K+1)$ term, including the displayed prefactor.

For completeness, let
\[
 A_0={}_1F_1(1/4;1/2;1/(4t)),\quad
 B_0={}_1F_1(3/4;3/2;1/(4t)),
\]
and let $C_0,D_0$ have parameters $(5/4;3/2)$ and $(3/4;1/2)$,
respectively, at the same argument. A rational-coefficient local frame is
\begin{equation}\label{eq:infinityframe}
 Y_\infty=\begin{pmatrix}
 t^{-1/4}A_0&t^{-3/4}B_0\\
 -\tfrac12t^{-5/4}C_0&-2t^{-3/4}D_0
 \end{pmatrix},\qquad
 \det Y_\infty=-2t^{-1}e^{1/(4t)}.
\end{equation}
The actual period frame is
\begin{equation}\label{eq:comparison}
 P=\frac1{\sqrt2}Y_\infty
 \begin{pmatrix}\Gamma(1/4)&i\Gamma(1/4)\\
 -\Gamma(3/4)&i\Gamma(3/4)\end{pmatrix}.
\end{equation}
In particular, $\det P=-2\sqrt2\pi i\,t^{-1}e^{1/(4t)}$.
The rational local frame and the contour-dependent comparison constants
have distinct roles.

The strong coefficients have a weighted combinatorial encoding, rather
than the same Wick interpretation. For example,
$4^n(1/4)_n=\prod_{j<n}(4j+1)$ counts histories of increasing plane
5-ary trees: insertion at one free leaf replaces it by five, increasing
the number of available leaves by four. The product
$\prod_{j<n}(4j+3)$ starts with three ordered leaves. The factorials and
powers in ${}_1F_1$ supply the weights. Equation \eqref{eq:Taylor}
likewise expands into finite weighted compositions of rational matrices.
These descriptions show exactly how different combinatorial series can
come from the same differential object.

\section{All finite Hermitian matrix sizes}\label{sec:matrix}
Let $\mathcal H_M$ be the real Euclidean space of Hermitian matrices,
with inner product $\Tr(XY)$. Its dimension is $d=M^2$. Define
\begin{equation}\label{eq:matrix}
 Z_M(\lambda)=
 \frac{\int_{\mathcal H_M}e^{-M\Tr(X^2)/2-M\lambda\Tr(X^4)/4}\dd X}
 {\int_{\mathcal H_M}e^{-M\Tr(X^2)/2}\dd X},\qquad\lambda>0.
\end{equation}
The eigenvalue determinant formula \cite[Section 5.1]{EKR}, followed by
$x\mapsto x/\sqrt M$, gives
\begin{equation}\label{eq:Hankel}
 Z_M(\lambda)=\frac{\det[q_{i+j}(\lambda/M)]_{0\le i,j<M}}
 {\prod_{j=0}^{M-1}j!}.
\end{equation}
Here odd moments vanish. The denominator is the Gaussian moment
determinant. By \eqref{eq:ibp}, every entry is a rational linear
combination of $q_0,q_2$. Thus $Z_M$ is a nonzero homogeneous polynomial
of degree $M$ in the same two periods, after the rational pullback
$t=\lambda/M$.

\begin{proposition}\label{prop:minimalODE}
The minimal homogeneous linear differential equation for $Z_M$ over
$\overline{\Q}(\lambda)$ has order exactly $M+1$.
The associated module is $\Sym^M$ of the rank-two quartic module.
Its formal exponential factors at zero are
$e^{j/(4t)}$, $0\le j\le M$.
\end{proposition}
\begin{proof}
Homogeneous monomials of degree $M$ in $q_0,q_2$ span a differential
module of rank $M+1$. The two periods are algebraically independent
over $K$ by Proposition~\ref{prop:galois}: after extending constants to
$\C$ the group is still $\GL_2$, and a first column of its generic
frame has no polynomial relation. Thus the indicated monomials
are linearly independent. The representation $\Sym^M$ of $\GL_2$ is
irreducible in characteristic zero. Consequently the differential
submodule generated by the nonzero polynomial \eqref{eq:Hankel} is the
whole module, so its cyclic equation has order $M+1$. The exponential
factors follow by taking the symmetric power of the two formal factors
at zero.
\end{proof}

At size two, for example, $Z_2(\lambda)=q_0(\lambda/2)q_2(\lambda/2)$,
and elimination gives
\begin{equation}\label{eq:M2ODE}
 \begin{split}
 8\lambda^4 Z_2'''+4\lambda^2(16\lambda+3)Z_2''
 +2(55\lambda^2+18\lambda+2)Z_2'
 +3(10\lambda+3)Z_2=0.
 \end{split}
\end{equation}
This has been checked symbolically against the symmetric-square system.
The determinant and holonomic methods are consistent with, and do not
replace, the established Painlev\'e description of quartic weights
\cite{ClarksonJordaan}.

Write $Z_M\sim\sum_{n\ge0}(-\lambda)^n w_{M,n}$.
With Gaussian covariance
$\E[X_{ij}X_{kl}]=M^{-1}\delta_{il}\delta_{jk}$, Wick's rule yields
\begin{equation}\label{eq:Wickmatrix}
 w_{M,n}=\frac1{4^n n!}\sum_{\alpha}M^{F(\sigma\alpha)-n}.
\end{equation}
Here $\sigma$ is a product of $n$ specified 4-cycles, $\alpha$ runs over
pairings of their $4n$ darts, and $F$ counts permutation cycles. The sum
includes disconnected diagrams. In particular,
$w_{M,0}=1$ and $w_{M,1}=(2M^2+1)/4$.
Connected diagrams instead describe $\log Z_M$.

This matters when choosing restrictions. Retaining a genus range changes
the coefficients in \eqref{eq:Wickmatrix}; it need not retain all moment
positivity conditions used below. Likewise, Tanasa's multi-matrix
schemes retain color, strand and insertion conditions
\cite{Tanasa,BNT}. Such conditions must be imposed before a graph sum
is identified with a particular integral. A formal interaction which
is not positive cannot be assigned the positivity bounds of this paper.

\section{Classification by positive kernel order}\label{sec:positive}
The next theorem gives a second precise meaning to the choice of a
restriction. It concerns a class of positive coefficient realizations
and resummations, not a literal base change of fields.

\begin{definition}
A normalized positive generalized Stieltjes representation of order
$\beta>0$ is
\begin{equation}\label{eq:Sbeta}
 f(\lambda)=\int_{[0,\infty)}(1+\lambda s)^{-\beta}\dd\rho_\beta(s),
 \qquad\rho_\beta([0,\infty))=1.
\end{equation}
Its exact order is the infimum of admissible $\beta$.
\end{definition}
This is the normalized convention for generalized Stieltjes functions
\cite{Sokal,KarpPrilepkina,KarpHypergeometric}. If all moments exist, its unsigned
asymptotic coefficients and representing moments satisfy
\begin{equation}\label{eq:reweight}
 f(\lambda)\sim\sum_{n\ge0}(-\lambda)^n w_n,
 \qquad \mu_{\beta,n}:=\int s^n\dd\rho_\beta(s)
                  =\frac{n!w_n}{(\beta)_n}.
\end{equation}
Thus a choice of order prescribes a particular reweighting of the
same combinatorial coefficients.

\begin{theorem}\label{thm:exactorder}
Let $Q$ be a strictly positive homogeneous polynomial of degree $2m$ on
$\mathbb R^d\setminus\{0\}$, with $m\ge2$, and let $X$ be a standard
Gaussian vector. For
$f_Q(\lambda)=\E e^{-\lambda Q(X)}$, the admissible orders in
\eqref{eq:Sbeta} are exactly
\[
 \beta\in[d/(2m),\infty).
\]
Put $a=d/(2m)$. A representing random variable at the minimal order is
\begin{equation}\label{eq:rho}
 T=(2m)^m Q(\Omega)\prod_{j=1}^{m-1}G_{a+j/m},
\end{equation}
where $\Omega$ is uniform on the unit sphere and the $G_b$ are
independent gamma variables of shape $b$ and scale one, independent of
$\Omega$. For every $\beta>a$, a representing variable is
$T B$, with independent $B\sim\operatorname{Beta}(a,\beta-a)$.
\end{theorem}
\begin{proof}
Write $X=R\Omega$. Then $S=R^2/2$ has the gamma law of shape $d/2=ma$,
independently of $\Omega$. Gamma multiplication gives the identity in law
\begin{equation}\label{eq:gammaproduct}
 S^m\ \stackrel{\mathrm{law}}=\ m^m\prod_{j=0}^{m-1}G_{a+j/m}.
\end{equation}
For rigor, this identity follows from the gamma multiplication formula
at imaginary Mellin arguments, hence from equality of the characteristic
functions of the logarithms. It does not assume that positive integer
moments determine the law. Since
$Q(X)=2^m S^m Q(\Omega)$, conditioning on the gamma variable $G_a$
proves $f_Q=\E(1+\lambda T)^{-a}$.

Euler's beta identity gives
\begin{equation}\label{eq:betainclusion}
 (1+\lambda s)^{-a}
   =\E_B(1+\lambda sB)^{-\beta},\qquad \beta>a,
\end{equation}
which proves inclusion in every larger order. This is the positive
probability form of the usual inclusion of Stieltjes classes
\cite[Theorem 3]{KarpPrilepkina}.

Finally, scaling $X=\lambda^{-1/(2m)}Y$ gives
\begin{equation}\label{eq:strongasymptotic}
 \lambda^a f_Q(\lambda)\longrightarrow
 c_Q=(2\pi)^{-d/2}\int_{\mathbb R^d}e^{-Q(Y)}\dd Y\in(0,\infty).
\end{equation}
Strict positivity on the sphere justifies dominated convergence and
finiteness. If a representation of order $\beta<a$ existed, monotone
convergence would give
$\lambda^\beta f_Q(\lambda)\to\int s^{-\beta}\dd\rho_\beta(s)>0$,
possibly infinite. This contradicts \eqref{eq:strongasymptotic}.
\end{proof}

\begin{corollary}\label{cor:matrixorder}
The exact order for \eqref{eq:matrix} is $M^2/4$. It is an ordinary
Stieltjes function if and only if $M=1$ or $M=2$. Its minimal-order
representing variable is
\begin{equation}\label{eq:matrixrho}
 T_M=\frac4M\Tr(\Omega^4)G_{(M^2+2)/4},
 \qquad\Tr(\Omega^2)=1.
\end{equation}
\end{corollary}
\begin{proof}
Set $Y=\sqrt M X$ in \eqref{eq:matrix} and apply the theorem to
$Q(Y)=\Tr(Y^4)/(4M)$ on its $M^2$ real coordinates. On the unit sphere,
$1/M\le\Tr(\Omega^4)\le1$.
\end{proof}

Already this classification excludes a tempting approximation argument:
at size three, ordinary Stieltjes--Pad\'e convergence cannot be invoked
for $Z_3$. The obstruction is its decay $\lambda^{-9/4}$, rather than
a failure of complete monotonicity. The generalized positive
representation remains available.

\section{Choosing the strongest positive moment restriction}\label{sec:best}
Fix an admissible order $\beta$ and suppose that the coefficients
$w_0,\ldots,w_{2n}$ are known. From \eqref{eq:reweight} construct the
$n$-node Gauss rule for $\rho_\beta$, denoted $G_{\beta,n}$, and the
$(n+1)$-node left Gauss--Radau rule $R_{\beta,n}$ with fixed node zero.
The latter's nonzero nodes are the Gauss nodes for $s\dd\rho_\beta(s)$.
Apply both rules to $h_\beta(s)=(1+\lambda s)^{-\beta}$ and write
$L_{\beta,n}=G_{\beta,n}(h_\beta)$ and
$U_{\beta,n}=R_{\beta,n}(h_\beta)$.
Assume the relevant Hankel moment matrices are positive definite; the
measures in Theorem~\ref{thm:exactorder} have this property.

\begin{lemma}\label{lem:extremal}
For every $\lambda>0$,
\[
 L_{\beta,n}\le f(\lambda)\le U_{\beta,n}.
\]
Moreover, $L_{\beta,n}$ is the smallest possible value of
$\int h_\beta\dd\nu$ over positive measures sharing moments through
degree $2n-1$, and $U_{\beta,n}$ is the largest possible value over
those sharing moments through degree $2n$.
\end{lemma}
\begin{proof}
Hermite interpolation at the Gauss nodes, with multiplicity two, has
degree at most $2n-1$. Its remainder has the sign of
$h_\beta^{(2n)}>0$ times the square of the nodal polynomial.
It therefore lies below $h_\beta$ on $[0,\infty)$.
Integrating against any measure with the prescribed moments gives
the lower bound. The discrete Gauss measure has those moments and
attains equality, proving the extremal assertion.

For Radau interpolation use the simple node zero and double nonzero
nodes. The remainder has the sign of
$s h_\beta^{(2n+1)}<0$ times a square, so the interpolant is an upper
bound. The discrete Radau measure matches moments through degree $2n$
and attains it. All integrals exist since the interpolants have the
specified degrees. This is the classical sign mechanism for positive
quadrature bounds \cite{Gautschi,LRW}.
\end{proof}

\begin{theorem}[Order selection]\label{thm:selection}
Suppose that $f$ has a representation of order $a>0$ with all moments
finite. For $a\le\beta_1<\beta_2$, form the moments at each order from
the same coefficients $w_j$ by \eqref{eq:reweight}. Then
\begin{equation}\label{eq:dominance}
 L_{\beta_2,n}\le L_{\beta_1,n}\le f(\lambda)
 \le U_{\beta_1,n}\le U_{\beta_2,n},\qquad\lambda>0.
\end{equation}
Hence the smallest admissible order gives the strongest bounds in
this hierarchy, at each truncation and each positive coupling.
\end{theorem}
\begin{proof}
For any positive measure $\nu_1$ with the prescribed order-$\beta_1$
moments through degree $k$, let $\nu_2$ be the law of $SB$ where
$S\sim\nu_1$ and
$B\sim\operatorname{Beta}(\beta_1,\beta_2-\beta_1)$ independently.
Its moments are
\[
 \int s^j\dd\nu_2=\frac{(\beta_1)_j}{(\beta_2)_j}
                   \frac{j!w_j}{(\beta_1)_j}
                 =\frac{j!w_j}{(\beta_2)_j},\qquad j\le k.
\]
The beta identity \eqref{eq:betainclusion} preserves the integral value.
Thus every value attainable in the order-$\beta_1$ moment class is
also attainable in the order-$\beta_2$ moment class. Taking infima
for $k=2n-1$ and suprema for $k=2n$, and applying
Lemma~\ref{lem:extremal}, proves \eqref{eq:dominance}.
\end{proof}

This is a uniform selection rule, with no estimated crossover parameter.
It says that forgetting the beta factorization enlarges the possible
values compatible with the coefficient data. At order one the rules
are the familiar Stieltjes Pad\'e bounds $[n-1/n]$ and $[n/n]$
\cite{Allen}. At nonintegral orders they are
positive sums of fractional powers, so they need not be rational
functions of $\lambda$.

\begin{theorem}\label{thm:convergence}
For the quartic and sextic cases of Theorem~\ref{thm:exactorder}, the
minimal-order bounds converge monotonically to $f_Q(\lambda)$ at every
$\lambda>0$. The same holds at every larger admissible order.
\end{theorem}
\begin{proof}
The random variable $Q(\Omega)$ is bounded. Formula \eqref{eq:rho}
therefore gives, for a fixed constant $C$, moments bounded by
\[
 \mu_{a,n}\le C^n\prod_{j=1}^{m-1}(a+j/m)_n.
\]
For $m=2,3$, Stirling's formula implies
$\mu_{a,n}^{-1/(2n)}\ge c\,n^{-(m-1)/2}$ for all sufficiently large
$n$. The Stieltjes Carleman series diverges, so the representing measure
is determinate. At higher order the beta factor is bounded by one and
does not worsen this estimate.

By the extremal description, increasing $n$ increases the lower bound
and decreases the upper bound. The discrete Gauss and Radau measures
are tight since their first moments are fixed. For any fixed moment,
a higher fixed even moment bounds the tails uniformly as $n$ grows;
thus every weak limit has all the prescribed moments. Determinacy
forces every such limit to be $\rho_\beta$. Finally, $h_\beta$ is bounded
and continuous, so both quadratures converge to its integral.
\end{proof}

\section{A sextic consequence and a limit of the method}\label{sec:sextic}
Let $X$ be a standard real Gaussian and define
\[
 P_{m,k}(\lambda)=\frac{\E[X^{2k}e^{-\lambda X^{2m}/(2m)}]}
 {(2k-1)!!},\qquad k\ge0,
\]
with $(-1)!!=1$. This normalization gives $P_{m,k}(0)=1$.
Under the $X^{2k}$-weighted Gaussian measure, $X^2/2$ has gamma
shape $k+1/2$. The proof of Theorem~\ref{thm:exactorder} applies
unchanged. Its exact order and representing variable are
\begin{equation}\label{eq:scalarhigher}
 a_k=\frac{2k+1}{2m},\qquad
 T_{m,k}=(2m)^{m-1}\prod_{j=1}^{m-1}G_{a_k+j/m}.
\end{equation}
All its moments come explicitly from the Wick coefficients
\begin{equation}\label{eq:scalarhigherwick}
 w_{m,k,n}=\frac{(2k+2mn-1)!!}
 {(2k-1)!!(2m)^n n!}.
\end{equation}

\begin{corollary}\label{cor:sextic}
For the scalar sextic model ($m=3$), every normalized even correlation
function admits convergent, certified bounds at all positive couplings:
\begin{equation}\label{eq:ratio}
 (2k-1)!!\frac{L_{a_k,n}^{(k)}}{U_{a_0,n}^{(0)}}
 \ \le\ \langle x^{2k}\rangle_\lambda
 \ \le\ (2k-1)!!\frac{U_{a_k,n}^{(k)}}{L_{a_0,n}^{(0)}}.
\end{equation}
The numerator and denominator rules use \eqref{eq:scalarhigherwick}.
Their endpoints converge to the correlation function.
\end{corollary}
\begin{proof}
Each factor is positive, and the same two-gamma Carleman argument in
Theorem~\ref{thm:convergence} applies for every fixed $k$.
Divide the respective bounds and use
$\langle x^{2k}\rangle=(2k-1)!!P_{3,k}/P_{3,0}$.
\end{proof}

The corollary supplies a concrete answer to part of the convergence
question suggested by the even-potential direction in \cite{FGL}.
Its approximants are generally algebraic functions, rather than the
rational Schwinger--Dyson approximants studied there. Convergence of
those particular sextic truncations, and their possible product
structure, remain separate questions.

There is also a precise reason not to extrapolate this argument to every
degree. For $m\ge4$, already the scalar minimal-order measure in
\eqref{eq:scalarhigher} is a product of at least three gamma variables.
It is moment indeterminate by the product criterion of Lin--Stoyanov
\cite[Corollary 1]{LinStoyanov}. Thus distinct positive measures share
all the reweighted Wick moments. The generalized Stieltjes transform
is injective on finite measures: substitute $\lambda=1/x$, divide the
result by $x^a$, and write $(x+s)^{-a}$ as a gamma Laplace integral.
Uniqueness of the two successive Laplace transforms proves the assertion.
These measures therefore give different functions for some positive
couplings, despite identical formal expansions.

This obstruction concerns reconstruction from that positive moment class.
It does not make the original integral ambiguous. The differential
equation, its analytic conditions, or the full gamma-product law
\eqref{eq:scalarhigher} can supply information absent from the integer
moments. In the proposed torsor approach, retaining such information is
essential: not every useful restriction is equally informative, even
when all its formal coefficients are known.

Exact relations also strengthen finite bounds in the determinate case.
For the sextic model, integration by parts gives
$\langle x^2\rangle+\lambda\langle x^6\rangle=1$.
An interval $[l_6,u_6]$ for the sixth moment therefore restricts the
second moment to $[1-\lambda u_6,1-\lambda l_6]$.
Intersecting this with its independently computed interval retains
information common to the moments. At $\lambda=10$, the two-node rules
give $0.0742<\langle x^2\rangle<1.4575$; the intersection improves the
upper bound to $0.9154$. The archived calculation certifies both
intervals. This elementary improvement uses an exact equation in addition
to the separate positive moment representations.

\section{Fields on rewrite graphs and hypergraphs}\label{sec:rewrite}
A rewrite system gives a directed graph whose vertices are states and
whose arrows are individual rule applications. Equal states are merged;
this convention matters when independent updates can be performed in
either order. Such state graphs, including those of hypergraph rewriting,
are studied in \cite[Sections 5.3 and 6.2]{Wolfram}. We use finite connected
portions and their underlying undirected graphs. The field model below
is an additional choice; the rewrite rules alone do not prescribe it.

Let $H$ be a connected graph with $n\ge2$ vertices and positive algebraic edge
conductances $c_e$. Fix a vertex $r$, put $x_r=0$, and set $d=n-1$ and
\[
 E_H(x)=\sum_{uv\in E(H)}c_{uv}(x_u-x_v)^2=x^{\mathsf T}L_rx,
 \qquad
 I_H(\lambda)=\int_{\mathbb R^d}
 e^{-E_H(x)/2-\lambda E_H(x)^2/4}\dd x.
\]
Here $L_r$ is the grounded weighted Laplacian. Write $\tau(H)$ for the
weighted spanning-tree sum, $R_H(u,v)$ for effective resistance, and
$\langle\cdot\rangle_{H,\lambda}$ for normalized expectation in this
integral. The quadratic model at $\lambda=0$ is a pinned Gaussian
field on a network \cite{LyonsPeres}.

\begin{proposition}\label{prop:graphfield}
Define $Z_d(\lambda)=\E\exp(-\lambda |Y|^4/4)$ for a standard Gaussian
$Y\in\mathbb R^d$. Then
\begin{align}
 I_H(\lambda)&=\frac{(2\pi)^{d/2}}{\sqrt{\tau(H)}}Z_d(\lambda),
 \label{eq:graphpartition}\\
 \langle(x_u-x_v)^2\rangle_{H,\lambda}
 &=\frac{Z_{d+2}(\lambda)}{Z_d(\lambda)}R_H(u,v).
 \label{eq:graphcorrelation}
\end{align}
The exact positive-kernel order of $Z_d$ is $a=d/4$, with
\begin{equation}\label{eq:graphmoments}
 Z_d(\lambda)=\E(1+4\lambda G_{a+1/2})^{-a},\qquad
 w_{d,j}=\frac{4^j(a)_j(a+1/2)_j}{j!}.
\end{equation}
If $n$ is even, then
\begin{equation}\label{eq:graphperiod}
 Z_d(\lambda)=\frac{q_{d-1}(\lambda)}{(d-2)!!}.
\end{equation}
Consequently $Z_d$ and the normalized even
polynomial correlations are rational expressions in $q_0,q_2$.
The same differential period torsor applies at every even vertex number.
\end{proposition}
\begin{proof}
The matrix-tree theorem gives $\det L_r=\tau(H)$; set $y=L_r^{1/2}x$.
Rotational symmetry gives
$\langle yy^{\mathsf T}\rangle=(Z_{d+2}/Z_d)I$.
The usual network identity
$(e_u-e_v)^{\mathsf T}L_r^{-1}(e_u-e_v)=R_H(u,v)$, with the rooted
coordinate omitted, proves \eqref{eq:graphcorrelation}.
Theorem~\ref{thm:exactorder} applied to $Q(y)=|y|^4/4$ proves
\eqref{eq:graphmoments}, and radial integration proves
\eqref{eq:graphperiod}. For odd $d$, recurrence \eqref{eq:ibp}
reduces every required even moment to $q_0,q_2$.
For higher correlations the spherical moment formula gives
$Z_{d+2k}/Z_d$ times the Gaussian Wick sum of products of entries
of $L_r^{-1}$. Finally, the adjacent moment pairs are related by
\[
 \binom{q_{2j+2}}{q_{2j+4}}
 =\begin{pmatrix}0&1\\(2j+1)/\lambda&-1/\lambda\end{pmatrix}
   \binom{q_{2j}}{q_{2j+2}}.
\]
The determinant is nonzero, so these pairs are rational changes of
basis of the original rank-two connection.
\end{proof}

The graph factors have their own combinatorial descriptions.
The resistance is the spanning two-forest sum separating $u,v$, divided
by $\tau(H)$ \cite[Theorem 2]{BarrettForests}. A killed-random-walk
expansion of $L_r^{-1}$ gives another convergent series
\cite[Chapter 2]{LyonsPeres}. Either can be combined with the period
bounds. If $L_d\le Z_d\le U_d$ and
$L_{d+2}\le Z_{d+2}\le U_{d+2}$ are positive enclosures, then
\[
 R_H(u,v)\frac{L_{d+2}}{U_d}
 \le\langle(x_u-x_v)^2\rangle_{H,\lambda}
 \le R_H(u,v)\frac{U_{d+2}}{L_d}.
\]

\subsection*{Independent and dependent rewrites}
Consider the string rule $A\to AB$. Starting from $A^k$, each state
has the unique form $AB^{m_1}\cdots AB^{m_k}$, and one update increases
one $m_i$. Retaining the states between $A^k$ and $(AB)^k$ gives the
oriented cube $Q_k$. It has $2^k$ states and $k!$ complete histories.
Equivalently, represent a word by a directed path with edge colors
$A,B$, and subdivide an $A$-edge into consecutive $A$- and $B$-edges.
Isomorphism of these directed colored paths is equality of words.
The examples are therefore also graph-rewrite systems with unnamed vertices.
For $k=2$ the two paths are
\[
 AA\longrightarrow ABA\longrightarrow ABAB,
 \qquad
 AA\longrightarrow AAB\longrightarrow ABAB.
\]
The underlying graph is $C_4$. In contrast,
$A\to AB\to ABB\to ABBB$ gives $P_4$, a four-state chain.
Both use $Z_3=q_2$, but their spanning-tree counts are $4$ and $1$,
and their initial-to-terminal resistances are $1$ and $3$.
Thus $I_{C_4}=I_{P_4}/2$, whereas the chain's endpoint variance is
three times the diamond's at every positive coupling. At $\lambda=1$,
16-node bounds give respectively
$0.3777<\langle(x_u-x_v)^2\rangle_{C_4,1}<0.3809$ and
$1.1333<\langle(x_u-x_v)^2\rangle_{P_4,1}<1.1425$.

Three independent rewrites already change the admissible approximation
class: $Q_3$ has $d=7$ and exact order $7/4$, so no ordinary Stieltjes
representation exists. There is a finite certificate. At ordinary
order one the candidate moments are $\mu_j=w_{7,j}$, and
\[
 \det(\mu_{i+j})_{0\le i,j\le2}
 =-\frac{3745374633}{65536}<0.
\]
At order $7/4$ all the required Hankel matrices are positive.
At coupling one, the 16-node interval has width below $0.000841$;
order $11/4$, using the same coefficients, has width above $0.00437$.
The example gives an exclusion test and an optimal choice among the
admissible orders.

For completeness, the cube's Laplacian eigenvalues are $2j$ with
multiplicity $\binom{k}{j}$, as follows by using the characters of
$(\mathbb Z/2)^k$. The matrix-tree theorem and the same characters give
\[
 \tau(Q_k)=2^{2^k-k-1}\prod_{j=1}^k j^{\binom{k}{j}},\qquad
 R_{Q_k}(u,\bar u)=\frac2{2^k}
       \sum_{\substack{1\le j\le k\\j\text{ odd}}}
          \frac{\binom{k}{j}}j.
\]
These formulas provide examples of arbitrary size within the same torsor.

\subsection*{Hypergraphs and the choice of interaction}
For a finite hypergraph $\mathcal H$, use its incidence graph
$B_{\mathcal H}$: vertices and hyperedges form its two vertex classes,
and incidence gives its edges. This construction retains the hyperedges
as separate objects. For
$V=\{0,1,2,3\}$ and $E=\{\{0,1,2\},\{1,2,3\}\}$, the unit-weight
incidence graph has six vertices, $\tau=4$, and $R(0,3)=3$.
Its partition function uses $Z_5=q_4/3=(q_0-q_2)/(3\lambda)$ and
has exact order $5/4$.

\begin{table}[ht]
\centering
\small\begin{tabular}{lrrrrr}
\toprule
$H$ & $\tau(H)$ & $R(u,v)$ & $a$ & lower $Z_d(1)$ & upper $Z_d(1)$\\
\midrule
$Q_1$ & 1 & $1$ & $1/4$ & 0.77198088 & 0.77212212\\
$Q_2=C_4$ & 4 & $1$ & $3/4$ & 0.36099863 & 0.36158743\\
$P_4$ & 1 & $3$ & $3/4$ & 0.36099863 & 0.36158743\\
$Q_3$ & 384 & $5/6$ & $7/4$ & 0.04460770 & 0.04544794\\
$Q_4$ & 42467328 & $2/3$ & $15/4$ & 0.00019404 & 0.00031685\\
$B_{\mathcal H}$ & 4 & $3$ & $5/4$ & 0.13660314 & 0.13747832\\
\bottomrule
\end{tabular}

\caption{Rewrite and incidence graphs with unit conductances. The
resistance is between initial and terminal states, between antipodes
for $Q_k$, or between vertices $0,3$ in the hypergraph example.
Bounds use 16 Gauss nodes and 17 Radau nodes at minimal order;
endpoints are rounded outward.}\label{tab:rewrite}
\end{table}

The square $E_H^2$ is a global interaction and has a precise limitation:
at fixed $n$, integrating out the field weights graphs by
$\tau(H)^{-1/2}$ independently of $\lambda$. The nonlinear coupling
changes field correlations but not these relative graph weights.
An edge-local quartic instead gives
\[
 \widetilde I_H(\lambda)=\int
 e^{-\frac12\sum_e y_e^2-\frac\lambda4\sum_e y_e^4}\dd x,
 \qquad y_e=\sqrt{c_e}(x_u-x_v).
\]
On a tree the $d$ edge differences are independent Gaussian coordinates,
so $\widetilde I_H=(2\pi)^{d/2}\tau(H)^{-1/2}q_0^d$.
This is a symmetric-power expression in the same periods, again with
exact order $d/4$. On graphs with cycles, this factorization generally
fails, although Theorems~\ref{thm:exactorder}--\ref{thm:convergence}
still apply. The accompanying example notes give the walk series,
leaf-history coefficients, and checks for the edge-local $P_4$ model.

\section{Certified computations and practical choices}\label{sec:computations}
The accompanying program constructs \eqref{eq:Hankel} with rational
truncated polynomials to obtain $w_{M,n}$, then forms
$\mu_{\beta,n}$ by \eqref{eq:reweight}. Exact rational Hankel systems
give monic orthogonal polynomials. Arb root isolation and positive
Christoffel weights enclose the quadrature nodes, weights and evaluations.
The program checks orthogonality and moment matching exactly or with
certified balls, as appropriate. No fitted integral values enter the
quadratures.

\begin{table}[ht]
\centering
\begin{tabular}{rrrrr}
\toprule
$M$ & $n$ & lower bound & upper bound & width $< $\\
\midrule
1 & 2 & 0.74828962 & 0.79361826 & $4.54\times10^{-2}$\\
1 & 4 & 0.76598846 & 0.77778969 & $1.19\times10^{-2}$\\
1 & 16 & 0.77198088 & 0.77212212 & $1.42\times10^{-4}$\\
2 & 2 & 0.37044526 & 0.45678743 & $8.64\times10^{-2}$\\
2 & 4 & 0.39920223 & 0.41762905 & $1.85\times10^{-2}$\\
2 & 16 & 0.40725400 & 0.40732727 & $7.33\times10^{-5}$\\
3 & 2 & 0.11288594 & 0.20355949 & $9.07\times10^{-2}$\\
3 & 4 & 0.13291925 & 0.15093682 & $1.81\times10^{-2}$\\
3 & 16 & 0.13889499 & 0.13893524 & $4.03\times10^{-5}$\\
4 & 2 & 0.02091036 & 0.08766673 & $6.68\times10^{-2}$\\
4 & 4 & 0.02819444 & 0.04032442 & $1.22\times10^{-2}$\\
4 & 16 & 0.03080806 & 0.03082700 & $1.90\times10^{-5}$\\
\bottomrule
\end{tabular}

\caption{Minimal-order bounds for $Z_M(1)$. The Gauss rule has $n$ nodes,
and Radau has $n+1$; moments through degree $2n$ suffice for the pair.
Displayed endpoints are rounded outward. Widths bound the unrounded
certified intervals.}\label{tab:quadrature}
\end{table}

The archived checks cover $1\le M\le4$, $1\le n\le16$,
$\lambda=1/100,1/10,1,10,100$, and orders
$a,a+1,a+2$, together with ordinary order one when it is admissible.
They verify the signs, monotonicity in $n$, and the order dominance in
Theorem~\ref{thm:selection}. Independent Bessel evaluations, inserted
in \eqref{eq:Hankel}, lie in every enclosure. The same program checks
the two independent formulas for sextic representing moments and
constructs bounds for the partition function and the second, fourth and
sixth correlations. Independently enclosed convergent strong expansions
lie in these bounds. These finite tests supplement
the proofs; they are not the convergence argument.

Positive moment bounds are not always the fastest evaluation. For the
quartic model we can compare certified Taylor bounds directly.
At size two let $w_n=w_{2,n}$ and define
\[
 B_n=\sum_{j=0}^n
 \frac{\Gamma(j/2+1/4)\Gamma((n-j)/2+3/4)}{j!(n-j)!}.
\]
The two expansions and their degree-$K$ absolute error bounds are
\begin{align*}
 Z_2(\lambda)&\sim\sum_{n\ge0}(-\lambda)^n w_n,
 &\epsilon_{\rm G}&\le\lambda^{K+1}w_{K+1},\\
 Z_2(\lambda)&=\frac1{\pi\lambda}
 \sum_{n\ge0}(-\sqrt{2/\lambda})^n B_n,
 &\epsilon_{\rm Q}&\le\frac{(2/\lambda)^{(K+1)/2}B_{K+1}}{\pi\lambda}.
\end{align*}
The second bound follows by expanding the quadratic part under the
positive product quartic measure, rather than by assuming that alternating
terms decrease. At an ordinary center $\lambda=1$, Cauchy's estimate on
$|\lambda-1|=R<1$ gives
\[
 \epsilon_{\rm T}\le
 \frac{(|\lambda-1|/R)^{K+1}}{1-|\lambda-1|/R},
 \qquad |\lambda-1|<R,
\]
because $|Z_M(\lambda)|\le1$ for $\Re\lambda\ge0$.
The Taylor coefficients use the seed periods at the center.

\begin{table}[ht]
\centering
\begin{tabular}{rrrr}
\toprule
$\lambda$ & Gaussian terms & quartic terms & Taylor terms at $1$\\
\midrule
$0.01$ & 10 & 591 & ---\\
$0.1$ & --- & 91 & ---\\
$0.9$ & --- & 29 & 11\\
$1$ & --- & 27 & 1\\
$1.1$ & --- & 26 & 11\\
$10$ & --- & 13 & ---\\
$100$ & --- & 8 & ---\\
\bottomrule
\end{tabular}
\caption{Terms sufficient to certify absolute error below $10^{-10}$
for $Z_2$, including Arb roundoff. Gaussian sums were searched through
degree 128, quartic sums through degree 1024. A dash denotes failure of
that bound in the search or exclusion from the chosen Taylor circle
$R=0.9$. Seed computation is excluded from the Taylor counts.}
\label{tab:charts}
\end{table}

For an actual choice, first fix the available coefficient data and the
cost to obtain new coefficients or seed periods. Within positive moment
methods use order $M^2/4$. Then compare its certified interval with
the weak, strong and ordinary Taylor bounds at the requested precision.
The narrowest interval at a fixed budget is a reproducible criterion;
it does not imply minimum running time without including coefficient,
root-isolation and seed costs. At $\lambda=1$, for instance, a short
convergent quartic expansion is much more accurate than a 16-node moment
rule. Moment rules instead give a seed-free method from Wick coefficients
which extends beyond the rank-two special-function example.
Rigorous special-function software already changes between expansions
and continuation methods \cite{Johansson}. Constructive and generalized
Borel methods provide other established reorganizations of perturbation
theory \cite{Rivasseau,Mera}. The existence of several quartic expansions
is not itself a new numerical method.

\section{Arithmetic restrictions and remaining questions}
The connection has an integral model over
$\mathbb Z[1/2,t,t^{-1}]$. It can therefore be reduced at any odd prime.
This is a change of base ring; a homomorphism $\Q\to\mathbb F_p$ does
not exist. Such a reduction need not admit the characteristic-zero local
frames.

\begin{proposition}\label{prop:pcurvature}
With $\nabla=D-A$, the trace of the reduced $p$-curvature is
\[
 \Tr\Psi_p=\frac1{4t^{2p}}\ne0\qquad(p\text{ odd}).
\]
Consequently the reduced connection has no full horizontal frame over
$\mathbb F_p((t))$ or $\mathbb F_p((t^{-1/4}))$ with their standard
derivations.
\end{proposition}
\begin{proof}
For the determinant connection write
$a=\Tr A=-t^{-1}-(4t^2)^{-1}$. Its $p$-curvature is
$-(D^{p-1}a+a^p)$. Wilson's identity gives
$D^{p-1}(-t^{-1})=t^{-p}$; the double-pole term differentiates to
zero. The simple-pole terms cancel with those of $a^p$, leaving the
displayed value. Determinants commute with the connection, so this is
the trace in rank two. In either displayed Laurent field $D^p=0$.
A full horizontal frame would trivialize the connection and force
its $p$-curvature to vanish, a contradiction.
\end{proof}

Finite rational jets can still be reduced whenever their denominators
are units. The companion computation verifies strong-basis jets through
degree 24 at $p=53,59,61$, and independently calculates the curvature
at $p=3,5,7,11$. Prime-field data can reconstruct finite rational
coefficients with denominator and height bounds; they do not directly
approximate a real exponential period. A Frobenius or exponential-sum
realization needs its own comparison theorem.

There are therefore two proved classifications here: the formal local
types of a specified differential torsor, and the admissible positive
kernel orders for a specified family of integrals. The second comes
with a selection theorem and, for quartic and sextic interactions,
convergence. Neither is a classification of every graph restriction or
every arithmetic realization.

The next questions are concrete. For constrained multi-matrix models,
can one compute the needed reweighted moments from finite schemes
while retaining positivity? Can extra algebraic information in the
period system sharpen the positive moment bounds, or remove the
indeterminacy obstruction at higher degree? For sextic models, how do
the certified quadratures compare with the particular rational
Schwinger--Dyson truncations, whose convergence is not proved here?
These questions concern further results, rather than consequences
already supplied by the word ``torsor.''

\bibliographystyle{abbrvurl}
\bibliography{integral-torsor}
\end{document}
